\documentclass[11pt]{article}
\usepackage[localtheorems]{raeez-math-template}
\usepackage{array}

\newtheorem{theorem}{Theorem}
\newtheorem{proposition}[theorem]{Proposition}
\newtheorem{lemma}[theorem]{Lemma}
\newtheorem{corollary}[theorem]{Corollary}
\newtheorem{conjecture}[theorem]{Conjecture}
\theoremstyle{definition}
\newtheorem{definition}[theorem]{Definition}
\newtheorem{remark}[theorem]{Remark}

\providecommand{\C}{\mathbb{C}}
\providecommand{\Z}{\mathbb{Z}}
\providecommand{\Q}{\mathbb{Q}}
\providecommand{\R}{\mathbb{R}}
\providecommand{\cA}{\mathcal{A}}
\providecommand{\cC}{\mathcal{C}}
\providecommand{\cH}{\mathcal{H}}
\providecommand{\cM}{\mathcal{M}}
\providecommand{\cO}{\mathcal{O}}
\providecommand{\DT}{\mathrm{DT}}
\providecommand{\CoHA}{\mathrm{CoHA}}
\providecommand{\Stab}{\mathrm{Stab}}
\providecommand{\Hilb}{\mathrm{Hilb}}
\providecommand{\Sp}{\mathrm{Sp}}
\providecommand{\sgn}{\mathrm{sgn}}
\providecommand{\fg}{\mathfrak{g}}
\providecommand{\kch}{\kappa_{\mathrm{ch}}}
\providecommand{\kcat}{\kappa_{\mathrm{cat}}}
\providecommand{\kBKM}{\kappa_{\mathrm{BKM}}}
\providecommand{\kfib}{\kappa_{\mathrm{fiber}}}
\providecommand{\Eone}{E_1}
\providecommand{\Etwo}{E_2}
\providecommand{\Ethree}{E_3}
\providecommand{\PhiFA}{\Phi^{\mathrm{FA}}}
\providecommand{\SpCh}{\mathrm{Sp}^{\mathrm{ch}}}

\title{DT Wall-Crossing and the Borcherds Denominator for $K3\times E$}
\author{Raeez Lorgat}
\date{\today}

\begin{document}
\maketitle

\begin{abstract}
The reduced Donaldson--Thomas partition function of $K3\times E$ is
compared with the Weyl--Kac--Borcherds denominator of
$\fg_{\Delta_5}$. The proved $(N,M)=(1,1)$ case is separated from the
CHL and Gritsenko--Clery extensions. The Kontsevich--Soibelman product
belongs to the Borcherds-product side of the denominator identity; the
Macdonald alternating sum is the Weyl-group character-formula avatar of
the same automorphic form. The four scalar invariants attached to
$K3\times E$ are
\[
 \{\kcat,\kch^{\mathrm{Heis}},\kBKM,\kfib\}=\{0,3,5,24\},
\]
with the fibre value $\kcat(K3)=2$ kept separate.
\end{abstract}

\section{Sources and Conventions}
\label{sec:sources-conventions}

The source hierarchy for the formulas below is:
Oberdieck--Pandharipande for reduced DT theory on $K3\times E$,
Oberdieck--Pixton for the holomorphic-anomaly proof of the Igusa cusp
form conjecture, Gritsenko--Nikulin and Borcherds for the
denominator identity, Gritsenko for the multiplicative lift, and
Kontsevich--Soibelman for motivic wall-crossing. The Igusa repository is
used only as a consistency check; it is not a source of proof.

Let $X=K3\times E$. The stage-one object
$\PhiFA_3(D^b\mathrm{Coh}(X))$ is a holomorphic factorisation algebra
on the Calabi--Yau threefold. A choice of $(\Sigma_2,C)$ gives the
specialised chiral output
\[
 \Phi_3^{(\Sigma_2,C)}(D^b\mathrm{Coh}(X))
 = \SpCh_{\Sigma_2,C}\bigl(\PhiFA_3(D^b\mathrm{Coh}(X))\bigr),
\]
an $\Eone$-chiral algebra on $C$. The $\Etwo$ structure attached to the
same data lives on the Drinfeld centre or double under the stated
pairing hypotheses; it is not native on the specialised $\Eone$ output.

\section{Reduced DT Partition Function}
\label{sec:reduced-dt}

\begin{definition}[Reduced DT zeta]
Let $\beta_h\in H_2(K3,\Z)$ be a primitive curve class with
$\langle \beta_h,\beta_h\rangle=2h-2$. The reduced DT partition
function is
\[
 Z^{X}_{\mathrm{red}}(q,t,p)
 =\sum_{h,d\ge 0}\sum_{n\in\Z}
 \DT^{\mathrm{red}}_{n,(\beta_h,d)}
 q^{d-1}t^{h-1}(-p)^n,
\]
where
\[
 \DT^{\mathrm{red}}_{n,(\beta_h,d)}
 =\int_{\Hilb^n(X,(\beta_h,d))/E}\nu_B
\]
is the Behrend-weighted Euler characteristic after quotienting the
translation action of $E$.
\end{definition}

\begin{theorem}[The primitive $K3\times E$ identity]
\label{thm:primitive-k3e-dt}
For $X=K3\times E$ and primitive K3 class, Oberdieck--Pandharipande
and Oberdieck--Pixton prove
\[
 Z^{X}_{\mathrm{red}}(q,t,p)
 = C_{\mathrm{OP}}\Delta_5(q,t,p)^{-2}
\]
for the normalisation constant $C_{\mathrm{OP}}$ fixed by the chosen
Igusa cusp form convention. In the monic product convention
$D_5=64^{-1}\Delta_5$ used in the Oberdieck--Pandharipande scalar,
\[
 \Phi_{10}^{\mathrm{OP}}=D_5^2=4096^{-1}\Delta_5^2,
 \qquad
 Z_{\mathrm{OP}}(K3\times E)=-(\Phi_{10}^{\mathrm{OP}})^{-1}
 =-4096\,\Delta_5^{-2}.
\]
The unreduced invariant is zero because the $E$-translation direction
contributes a trivial obstruction factor.
\end{theorem}

\section{Product, Alternating Sum, and Wall-Crossing}
\label{sec:product-sum-wall}

Let $\Lambda^{2,1}_{II}$ be the hyperbolic rank-three lattice in the
Gritsenko--Nikulin construction. The denominator identity for
$\fg_{\Delta_5}$ has two expressions:
\[
 \Delta_5(Z)
 = e^{-2\pi i(\rho,Z)}
 \prod_{\alpha\in\Delta^+}
 \left(1-e^{-2\pi i(\alpha,Z)}\right)^{\mathrm{mult}(\alpha)}
\]
and
\[
 \Delta_5(Z)
 = \sum_{w\in W^{(2)}(\Lambda^{2,1}_{II})}\sgn(w)
 \left(e^{-2\pi i(w\rho,Z)}
 -\sum_{a}m(a)e^{-2\pi i(w(\rho+a),Z)}\right).
\]
The product is the Borcherds expansion at a chamber. The alternating sum
is the Weyl--Kac--Borcherds character formula after summing over Weyl
chambers.

\begin{proposition}[Where KS wall-crossing enters]
\label{prop:ks-product-side}
Kontsevich--Soibelman wall-crossing acts on the ordered product of
automorphisms attached to active DT charges. Under the
Gritsenko--Nikulin identification at the primitive $K3\times E$ slice,
this is the Borcherds-product side of the denominator identity. The
Macdonald alternating sum enters after passing from one chamber product
to the Weyl-group orbit of chamber products.
\end{proposition}

\begin{proof}
The KS formula preserves the ordered product
\[
 \prod_{\ell\subset V}^{\curvearrowleft}
 \exp\!\left(\sum_{Z(\gamma)\in \ell}\Omega(\gamma)e_\gamma\right)
\]
when the stability condition crosses a wall. This is a reordering
statement in the completed lattice Lie algebra. The Borcherds product is
also an ordered chamber product, with exponents equal to root
multiplicities. The alternating sum is a separate consequence of the
Weyl--Kac--Borcherds denominator identity for $\fg_{\Delta_5}$.
\end{proof}

\section{The Wall-Crossing Gauge}
\label{sec:gauge}

\begin{definition}[DT lattice Lie algebra]
Let $\Gamma_X$ be the numerical charge lattice with Euler pairing
$\langle-,-\rangle_\chi$. The completed DT lattice Lie algebra is
\[
 \fg_{\Gamma_X}
 =\widehat{\bigoplus_{\gamma\ne 0}\Q e_\gamma},
 \qquad
 [e_{\gamma_1},e_{\gamma_2}]
 =(-1)^{\langle\gamma_1,\gamma_2\rangle_\chi}
 \langle\gamma_1,\gamma_2\rangle_\chi e_{\gamma_1+\gamma_2}.
\]
\end{definition}

\begin{proposition}[Single-wall gauge, with hypotheses]
\label{prop:single-wall-gauge}
Assume:
\begin{enumerate}[label=\textup{(H\arabic*)}]
\item Bridgeland stability conditions $\sigma_0,\sigma_1$ on $D^b(X)$
lie in adjacent chambers.
\item The common wall has a strict convex active cone $C_W\subset\Gamma_X$.
\item The KS integration map from the motivic Hall algebra to
$\exp(\fg_{\Gamma_X})$ is defined on the completed cone.
\item The primitive class comparison with the Gritsenko--Nikulin lattice
identifies the relevant charges with roots of $\fg_{\Delta_5}$.
\end{enumerate}
Then the gauge carrying the Maurer--Cartan element from $\sigma_0$ to
$\sigma_1$ is
\[
 \lambda_W
 = \log\left(
 \prod_{\gamma\in C_W}^{\curvearrowleft}
 \exp\bigl(\Omega_{\sigma_0}(\gamma)e_\gamma\bigr)
 \right)
 \in \widehat{\fg_{\Gamma_X}}.
\]
Its graded character under the primitive lattice comparison is the
local logarithm of the corresponding Borcherds product,
\[
 \mathrm{ch}\bigl(\iota(\lambda_W)\bigr)
 = \log\!\left(
 \prod_{(n,l,m)\in C_W}
 \left(1-e^{2\pi i(nz_1+lz_2+mz_3)}\right)^{f(nm,l)}
 \right).
\]
\end{proposition}

\begin{remark}
This is a local single-wall statement. A global identification of all
Humbert divisors with all Bridgeland walls for $D^b(K3\times E)$ is a
separate theorem to prove. The local formula records the correct
object: a completed Lie-algebra gauge with a Borcherds-product
character, rather than a scalar logarithm asserted as an element.
\end{remark}

\section{Borcherds Weights and the Gritsenko--Clery Catalogue}
\label{sec:weights}

\begin{theorem}[CHL Borcherds weight]
\label{thm:chl-weight}
For the CHL frame shapes with $N\in\{1,2,3,4,6\}$, the Borcherds
weight is
\[
 \kBKM(\Phi_N)=\frac{c_N(0)}{2},
\]
where $c_N(0)$ is the constant term of the chosen Borcherds input
Jacobi form in the Gritsenko--Nikulin normalisation. At $N=1$,
$c_1(0)=10$, so
\[
 \kBKM(\Phi_1)=5=\mathrm{wt}(\Delta_5).
\]
\end{theorem}

The Gritsenko--Clery diagonal-divisor catalogue is a different
classification surface. Its eight weights are
\[
 (5,2,3,1,2,\tfrac12,\tfrac32,1).
\]
Only the common row $F_{1,1}=\Delta_5$ is identified without extra
data. A row map from the eight diagonal-divisor forms to twisted-twined
K3 elliptic genera must specify the pair $(g,h)$, the input Jacobi
form, its constant term, and the paramodular group. The equality
$\kBKM=c(0)/2$ then applies row by row after those data are supplied.

\section{Conjecture for Twisted-Twined Quotients}
\label{sec:conjecture}

\begin{conjecture}[Twisted-twined DT denominator]
\label{conj:twisted-twined-dt}
Let $S$ be a projective K3 surface with commuting finite symplectic
automorphisms $g$ and $h$, and let $E$ be an elliptic curve with the
corresponding translation data. Suppose the quotient
\[
 X_{g,h}=(S\times E)/\langle g\times\tau,\ h\times\mathrm{id}_E\rangle
\]
is a smooth Calabi--Yau threefold or admits a crepant reduced-DT theory.
Let $\phi_{g,h}^{K3}(\tau,z)=\sum c_{g,h}(n,l)q^n\zeta^l$ be the
twisted-twined K3 elliptic genus, and suppose its Borcherds lift is a
diagonal-divisor paramodular form $\Delta_{g,h}$ of weight
$c_{g,h}(0)/2$. Then the reduced DT zeta is expected to satisfy
\[
 Z^{X_{g,h}}_{\mathrm{red}}(q,t,p)
 = C_{g,h}\Delta_{g,h}(q,t,p)^{-2},
\]
and the positive-half Hall algebra attached to the same reduced theory
has root multiplicities governed by the Fourier coefficients of
$\phi_{g,h}^{K3}$.
\end{conjecture}

At $(g,h)=(1,1)$ this is Theorem~\ref{thm:primitive-k3e-dt}. For the
other rows, the missing input is not a value of the weight formula but
the construction of the reduced DT theory, the row map, and the
Hall--Borcherds positive half.

\section{Equivariance}
\label{sec:equivariance}

\begin{proposition}[$\Sp_4(\Z)$ transformation law at $(1,1)$]
\label{prop:sp4-equivariance}
Let $G=\begin{psmallmatrix}A&B\\ C&D\end{psmallmatrix}\in\Sp_4(\Z)$.
If
\[
 \Delta_5(G\cdot Z)=\nu_{\Delta_5}(G)\det(CZ+D)^5\Delta_5(Z),
\]
then the reduced DT scalar in Theorem~\ref{thm:primitive-k3e-dt}
transforms as
\[
 Z^X_{\mathrm{red}}(G\cdot Z)
 =\nu_{\Delta_5}(G)^{-2}\det(CZ+D)^{-10}
 Z^X_{\mathrm{red}}(Z).
\]
\end{proposition}

\begin{proof}
This is immediate from $Z^X_{\mathrm{red}}=C_{\mathrm{OP}}\Delta_5^{-2}$.
The sign and exponent are forced by inversion and squaring.
\end{proof}

\section{The Four Scalars on $K3\times E$}
\label{sec:four-scalars}

\[
\begin{array}{c|c|l}
\text{invariant} & \text{value} & \text{source}\\
\hline
\kcat(K3\times E) & 0 &
\chi(\cO_{K3})\chi(\cO_E)=2\cdot 0\\
\kch^{\mathrm{Heis}}(K3\times E) & 3 &
\text{Heisenberg--Mukai specialisation}\\
\kBKM(\fg_{\Delta_5}) & 5 &
c_1(0)/2=10/2=\mathrm{wt}(\Delta_5)\\
\kfib(K3) & 24 &
\mathrm{rank}\,H^\ast(K3,\Z)
\end{array}
\]

The raw compact Hodge supertrace on $K3\times E$ is zero:
\[
 \sum_q(-1)^qh^{0,q}(K3\times E)=1-1+1-1=0.
\]
The value $3$ is the Heisenberg--Mukai specialised chiral scalar, and
the value $2$ is the separate fibre datum
$\kcat(K3)=\chi(\cO_{K3})$. The additive formula
$\kBKM=\kch+\chi(\cO_{\mathrm{fiber}})$ is not valid; the Borcherds
weight is $c(0)/2$ in the chosen automorphic input.

\section{Proof Obligations}
\label{sec:proof-obligations}

\begin{enumerate}
\item Construct the global Bridgeland wall stratification on
$D^b(K3\times E)$ and prove its comparison with the Humbert stratification
of $\mathbb{H}_2$ beyond the primitive chamber.
\item Construct the compact $K3\times E$ Hall positive half, its Hopf
pairing, and its Hall--Drinfeld double. The Mukai self-mirror Yangian
is related evidence, not this object.
\item Extend the DT/Borcherds coefficient comparison from primitive K3
classes to imprimitive classes with the multiple-cover corrections
included.
\item Give the row map from the Gritsenko--Clery eight forms to
twisted-twined K3 elliptic genera before applying
$\kBKM=c(0)/2$ outside the common $\Delta_5$ row.
\end{enumerate}

\end{document}
